\documentclass[11pt]{article}
\usepackage[a4paper,margin=1in]{geometry}
\usepackage{microtype}
\usepackage{lmodern}
\usepackage{amsmath,amssymb}
\usepackage{graphicx}
\usepackage{booktabs}
\usepackage{siunitx}
\usepackage{hyperref}
\usepackage[round,authoryear]{natbib}
\usepackage{xurl}
\usepackage{caption}

\hypersetup{colorlinks=true,linkcolor=black,citecolor=black,urlcolor=blue}
\bibliographystyle{plainnat}

\title{Pre-Registered Cross-Regional Earthquake Precursor Evaluation:\\
Five Execution-Time Failure Modes, and Cross-Regional Confirmation of\\
Foreshock-Window Elevated Benioff Strain at\\
Macro AUC$\,=\,$0.78 / $z_\mathrm{block}\,=\,$8.4}
\author{Matthew Loftus\\
\small Cedar Loop LLC, Bremerton, WA, USA\\
\small \texttt{matthew.a.loftus@gmail.com}}
\date{\today}

\begin{document}
\maketitle

\begin{abstract}
We pre-registered a CSEP-style cross-regional earthquake-precursor
evaluation across six training and two test regions, six features, and
three null models, and ran it under the locked protocol with one
documented post-execution amendment. Execution surfaced five concrete
methodological failure modes catalogued for re-use: an inherited
overlap-rejection rule that collides with itself in event-dense
subduction zones; ANSS ComCat as a teleseismic catalog for non-US
regions; the Aki MLE $b$-value $N{\ge}30$ floor combined with 30-day
windows at moderate~$M_c$; structural inaccessibility of both
pre-registered test regions; and the pre-registered precursor window
$[t-30, t)$ admitting the immediate-foreshock period. The pre-registered
per-window scalar and joint-classifier panels yield clean nulls. A
single trivial scalar — log$_{10}$ Benioff energy in the final 5 days
of the 30-day pre-event window — achieves macro AUC $= 0.779$
($\mathrm{CI}_{95}\,[0.70, 0.86]$, $z_\mathrm{iid} = +11.0$,
$z_\mathrm{block} = +8.4$ under a circular-shift dependence-aware null,
paired $\Delta\mathrm{AUC}$ vs symmetric placebo $= +0.27$,
$\mathrm{CI}_{95}\,[+0.20, +0.32]$) across four LORO splits. A
TLS-style template-correlation scan over the same Benioff trajectory
\emph{underperforms} the trivial scalar by 0.075 AUC, documented as a
methodological ablation. Per-target-magnitude stratification under the
trivial scalar shows monotonic scaling: macro AUC $= 0.73$ for
$M{\in}[4.5, 5.0)$, $0.88$ for $M{\ge}5.0$. The headline finding is a
quantitative cross-regional confirmation of the foreshock-pervasiveness
phenomenon documented per-region by \citet{trugman2019}, with the
simplest possible scalar and a publicly verifiable pre-registration
chain of custody.
\end{abstract}

% =============================================================================
\section{Introduction}
\label{sec:intro}

The earthquake precursor literature has been historically constrained by
two structural problems: most claims are per-region (no cross-regional
generalization test) and post-hoc (no committed-in-advance protocol).
The Collaboratory for the Study of Earthquake Predictability
\citep{schorlemmer2007, mizrahi2024} solved pre-registration for
\emph{rate}-based forecasts; recent ionospheric-precursor work
\citep{tao2025iono} ported the discipline to a different physical
domain. No equivalent exists for \emph{feature}-based seismic precursor
work.

In parallel, the foreshock literature has converged on the empirical
finding that elevated seismicity precedes most moderate mainshocks.
\citet{trugman2019} used a 1.81-million-event high-resolution
template-matched catalog to show that 72\% of $M{\ge}4$ mainshocks in
southern California have foreshock activity statistically elevated
above local background. The phenomenon is theoretically anchored by
the ETAS cascade-triggering framework
\citep{helmstetter2003a, helmstetter2003b}, which derives expected
foreshock-rate properties from a single self-similar triggering process
without invoking a distinct precursor mechanism. Recent work continues
to refine the methodology \citep{lippiello2025, khan2025, moutote2021}.

We bridge these two threads. We pre-registered a CSEP-style six-feature
cross-regional evaluation on declustered $M{\ge}4.5$ targets across six
training regions (California, Cascadia, Japan, Chile, Turkey, Italy)
and two held-out test regions (Mexico, Alaska); the locked protocol
\citep{loftus2026prereg} commits feature definitions, region polygons,
station list, null models, magnitude-of-completeness procedure,
evaluation metric, and significance thresholds before any AUC was
computed. One amendment \citep{loftus2026pra2} re-ran the evaluation
after v1 surfaced two execution-time failure modes; the amendment was
committed before re-evaluation, but \emph{after} v1 had been executed,
so we describe it as a \emph{post-v1 amendment} rather than as binding
pre-registration in the strictest sense (\S\ref{sec:disc-prereg}). The
full chain of custody is publicly available as a single repository with
one commit per work session.

This paper reports two contributions. First, a five-mode catalogue of
how pre-registered cross-regional precursor evaluation breaks in
practice; each mode is documented with diagnostic evidence and
attempted mitigation. Second, a quantitative cross-regional confirmation
of the \citet{trugman2019} foreshock-pervasiveness finding, recovered
by the simplest scalar that captures the immediate-foreshock energy
release: log$_{10}$ Benioff energy in the final 5 days of the 30-day
pre-event window. We frame the result as confirmation rather than
discovery: the foreshock phenomenon is established, but the
cross-regional AUC, dependence-aware significance, and explicit
pre-registration chain of custody are novel.

% =============================================================================
\section{Pre-Registration Protocol}
\label{sec:prereg}

The locked v1 protocol \citep{loftus2026prereg} commits: ANSS ComCat
catalog 2000--2024 ingested at $M{\ge}1.5$; per-region $M_c$ via
Wiemer-Wyss \citep{wiemer2000} max-curvature $+$ Woessner-Wiemer
\citep{woessner2005} $+0.20$ correction $+$ explicit $b$-vs-$M_c$
plateau check; ZBZ 2013 nearest-neighbor declustering
\citep{zaliapin2013} with $b{=}1.0$, $d_f{=}1.6$,
$\log_{10}\eta_0{=}{-}5.0$; declustering applied to both the catalog
and the target set; six features (b-value via \citealp{aki1965} MLE
with \citealp{shibolt1982} standard error, b-drift, Benioff total $+$
curvature, $n$ above $M_c$, plus three waveform-derived: spectral
slope, whitened-spectrogram entropy, HHT IMF1 instantaneous frequency);
three null types (A: aftershock-free random; B: minimal exclusion;
C: colored-noise synthetic for waveform features only); macro-AUC across
qualifying training regions with cross-region bootstrap CI95,
within-region permutation $z$, and Bonferroni-corrected $\alpha$. Pre-reg
v1 also commits an explicit ``will not'' list (no $M_c$ tuning, no
station/region swaps, no feature additions, no window changes).

Amendment v2 \citep{loftus2026pra2}, committed before re-evaluation but
after v1 was executed, makes three changes grounded in v1 evidence:
ISC global bulletin replaces ANSS ComCat (fixing failure mode~\#2
below); the overlap-rejection forbidden zone shrinks from
$[t'-30, t'+60]$ to $[t', t'+60]$ (partial fix for failure mode~\#1
below); a per-region minimum of $N{\ge}8$ kept precursor windows is
required for inclusion in the macro pool. The chain is publicly
verifiable: v1 SHA \texttt{a4f1c6f}, amendment SHA \texttt{05a4b0f},
all subsequent commits cite both. We discuss the implications of the
post-v1 amendment in \S\ref{sec:disc-prereg}.

% =============================================================================
\section{Methods}
\label{sec:methods}

\subsection{Catalog and station selection}
\label{sec:methods-catalog}

We use the ISC global bulletin via \texttt{obspy.clients.fdsn} for all
six training and two test regions over 2000-01-01 to 2024-12-31.
Catalog ingestion at $M{\ge}1.5$ as pre-registered would have been
intractably large at ISC; we use $M{\ge}2.5$ as a declared deviation,
an empirical compromise between bulletin-pull cost and $M_c$
identifiability. Per-region $M_c$ is the chosen value where the
$b$-vs-$M_c$ curve plateaus to $|\Delta b| < 0.02$ between consecutive
$M_c$; values are recorded as artifacts and not retuned. Per-region
waveform routing was not pre-specified: we use NCEDC for California
(BK network, Berkeley-hosted), IRIS for Cascadia (UW.LON), Turkey
(IU.ANTO), and Italy. Italy's pre-registered station IV.AQU is gone
post-2009; we use MN.AQU on IRIS, the post-2009 MedNet successor at
the same physical L'Aquila Observatory, documented as 95\%-availability
fallback per pre-reg \S 3.

\subsection{Per-window scalar features}
\label{sec:methods-features}

All five catalog-derived scalars and three waveform-derived scalars are
computed per 30-day pre-event window per pre-reg \S 4. The waveform
features are aggregated across six 10-minute snapshots evenly spaced
across the window using the region's primary broadband Z channel,
preprocessed (instrument-response-removal to ground velocity, bandpass
0.5--12~Hz clamped to $f_s/3$ to dodge any anti-alias artifacts).
Per-snapshot Welch PSD log-log slope \citep{brune1970}, whitened-
spectrogram Shannon entropy \citep{bensen2007}, and EMD IMF1
instantaneous frequency \citep{huang1998} are aggregated to per-window
scalars by mean (slope, entropy) and median (IMF1 IF). Repeating-event
rate, per pre-reg \S 4.1, is excluded from the headline panel as
deferred infrastructure.

\subsection{Joint multi-feature classifier}
\label{sec:methods-clf}

In addition to per-feature scalar AUC, we train a joint classifier over
the six computable features (the $b$-pair excluded due to failure
mode~\#3 below) under leave-one-region-out (LORO) cross-validation.
Each region's features are z-scored within region using the union of
precursor and Null A windows in that region (per-region domain
adaptation); pooled samples from three training regions are used to
fit a logistic regression (\texttt{class\_weight="balanced"}) and a
200-tree random forest, then scored on the held-out region.

\subsection{Trivial Benioff scalar (headline method)}
\label{sec:methods-scalar}

For each window, sub-divide into six 5-day sub-windows; compute Benioff
energy $\sum \sqrt{E_i}$ in each sub-window from declustered events
$\ge M_c$, where $E_i = 10^{1.5 M_i + 4.8}$~J. The per-window scalar
is $\log_{10}$ of the Benioff sum in the last sub-window (days 25--30
of the 30-day window). This is intentionally minimal: a single scalar,
no template, no LORO-dependent fit. ROC-AUC is computed precursor vs
Null~A within each held-out region; the macro is the mean across the
four LORO splits.

\subsection{TLS-style template-correlation scan (ablation)}
\label{sec:methods-tls}

Pre-reg \S 3.2 promised a TLS-style matched filter applied to feature
time series \citep[cf. matrix-profile-based seismic template matching;]
[]{senobari2024}. We realize this via per-window 6-point trajectories
at 5-day sub-window resolution. Templates are LORO precursor-mean
trajectories (averaged across the three training regions); each held-
out window is scored by Pearson correlation with the template; ROC-AUC
is computed precursor vs Null~A. We report this as an ablation rather
than the headline because it is uniformly outperformed by the trivial
scalar of \S\ref{sec:methods-scalar} (\S\ref{sec:results-ablation}).

\subsection{Significance gates and dependence-aware null}
\label{sec:methods-gates}

Every macro statistic is evaluated against three concurrent gates per
pre-reg \S 6.3: cross-region bootstrap CI95-lower exceeds 0.5;
$|AUC - 0.5|$ exceeds 0.05; within-region permutation $z$-score
exceeds 3 in magnitude. Bonferroni correction is applied per pre-reg
\S 6.4. Following peer-review feedback, in addition to the
pre-registered i.i.d.\ label-permutation null, we report a
\emph{dependence-aware} null: within each region, sort windows by
$t_\mathrm{start}$, then circularly shift labels by a uniformly-random
offset before recomputing AUC. This preserves the temporal
autocorrelation of the per-region window sequence (so that nearby
windows that share local seismic context remain near each other under
the null) while breaking the precursor-mainshock association. The
resulting $z_\mathrm{block}$ is a more conservative significance
estimate than the i.i.d.\ $z_\mathrm{iid}$. We also report the paired
$\Delta\mathrm{AUC}$ between the headline scalar and a 30-day-shifted
control on shared bootstrap resamples
(\S\ref{sec:results-controls}).

% =============================================================================
\section{Results}
\label{sec:results}

\subsection{Five-mode catalogue of execution-time failures}
\label{sec:results-failure-modes}

\paragraph{Three operational lessons (modes \#1, \#2, \#4).} Three of
the five modes are operational gotchas familiar to anyone running
CSEP-style work in subduction zones, but their concrete combination
under our protocol made them load-bearing. Mode~\#1: under v1's
$[t'-30, t'+60]$ overlap-rejection rule, regions with moderate-to-high
target density lose nearly all precursor windows: Japan (4{,}563
declustered $M{\ge}4.5$ targets) loses all of them; Chile (3{,}275)
loses all; Cascadia (277) keeps only 4; Turkey (389) keeps 2.
Amendment 2 shrinks the zone to $[t', t'+60]$ and recovers Cascadia to
36 and Turkey to 30 kept precursor windows; Japan and Chile remain
inaccessible (1 each). Mode~\#2: under v1's ANSS ComCat catalog, Japan
$M_c = 4.55$ and Chile $M_c = 4.45$, both at the $M{\ge}4.5$ target
threshold; b-value features become structurally uncomputable.
Amendment 1 switches to ISC and recovers Japan/Chile $M_c$ to 3.5;
California $M_c$ rises slightly (2.75 $\to$ 3.50, an ISC US-coverage
gap relative to NCEDC) but California still qualifies. Mode~\#4: both
pre-registered test regions, Mexico (843 declustered targets) and
Alaska (847), reduce to one kept precursor window each under the
PRA-2 overlap rule, far below the $N{\ge}8$ minimum; this is mode~\#1
generalized to held-out regions, structurally inaccessible without
further amendments addressing subduction-zone seismicity density.

\paragraph{Failure mode \#3: $b$-feature $\times$ $M_c$-window
incompatibility.} At PRA-2 $M_c$ values combined with 30-day windows,
the median $n$ of events $\ge M_c$ is 4--5 in California, Cascadia,
and Italy, and 76 in Turkey. The $b$-value MLE's $N{\ge}30$ floor is
therefore satisfied in only 3\% of windows for the first three regions
and 73\% in Turkey. The \citet{gulia2019} $b$-drop signal does not
transfer to short windows at moderate $M_c$. The $b$-vs-Null A macro
AUC of 0.722 reported elsewhere in our panel is a degenerate-sample
artifact (Cascadia AUC = 1.00 computed on 1 precursor $\times$ 2 null
finite values). Researchers planning to test $b$-drop precursors at
moderate $M_c$ in short windows should expect this incompatibility.

\paragraph{Failure mode \#5: precursor-window-includes-foreshock-period.}
The pre-reg defines the precursor window as $[t-30, t)$, including the
foreshock period. Any signal detected in this window may therefore
reflect foreshock-period seismicity rather than long-distance precursor
in the strict 25-day-ahead sense. We address this with three controls
(\S\ref{sec:results-controls}); the result is unambiguous: under the
pre-reg's literal definition, our positive finding is a real signal,
but operationally it is foreshock detection. Future pre-registered work
should commit explicitly on whether a ``precursor window'' includes the
final 5--7 days, since the answer determines whether the test is for
foreshock detection or long-distance precursor detection.

\subsection{Per-window scalar and joint-classifier results}
\label{sec:results-scalar}

Across the eight-feature $\times$ two-null panel evaluated on the four
qualifying training regions (California, Cascadia, Turkey, Italy), zero
of sixteen tests pass 3$\sigma$ or Bonferroni at $\alpha = 0.05/128$.
Computable catalog features (Benioff total per pre-reg, Benioff
curvature, $n$ above $M_c$) lie within 0.05 of 0.5; waveform features
(slope, entropy, IMF1 IF) lie 0.45--0.49 vs Null A with consistent
direction across regions but no individual feature reaching 3$\sigma$.
The $b$ and $b$-drift features are uncomputable (failure mode \#3) and
their AUCs are degenerate.

The joint LR + RF classifier under LORO yields macro AUC $= 0.526$ (LR)
and $0.480$ (RF), with permutation $z = +0.97$ and $-0.71$
respectively; neither passes 3$\sigma$ or Bonferroni. The combination
of features carries no signal that the per-feature scalars miss.

\subsection{Trivial-scalar headline finding}
\label{sec:results-scalar-finding}

The trivial scalar of \S\ref{sec:methods-scalar} ($\log_{10}$ Benioff
in the last 5 days of the 30-day window) yields:

\begin{center}
\begin{tabular}{lcccc}
\toprule
Quantity & California & Cascadia & Turkey & Italy \\
\midrule
LORO AUC               & 0.811 & 0.905 & 0.691 & 0.710 \\
$n$ precursor          & 37    & 36    & 30    & 36    \\
$n$ Null A             & 200   & 200   & 200   & 200   \\
\bottomrule
\end{tabular}
\end{center}

Macro AUC = $0.779$, cross-region bootstrap $\mathrm{CI}_{95} = [0.701,
0.858]$. Significance: i.i.d.\ label-permutation $z_\mathrm{iid} =
+11.0$ ($p \approx 0$); circular-shift dependence-aware
$z_\mathrm{block} = +8.4$ ($p \approx 0$). The dependence inflation
amounts to $\sim$25\% in $z$, not catastrophic. The result passes the
3$\sigma$ pre-reg gate and Bonferroni at $\alpha = 0.05/8 = 6.25
\times 10^{-3}$ under either null.

\subsection{Window-shift, placebo, and foreshock-mask controls}
\label{sec:results-controls}

We applied three controls to the trivial-scalar finding. (i)~\emph{
Precursor shift back 30 days}: precursor windows shifted to
$[t-60, t-30)$, null windows untouched. Macro AUC collapses from 0.779
to 0.513; the signal does not generalize to a 30-day-shifted window.
(ii)~\emph{Symmetric placebo shift}: \emph{both} precursor and null
windows shifted back 30 days. Macro AUC = 0.499 (within 0.001 of
chance); this rules out a temporal-confound interpretation of (i).
(iii)~\emph{Foreshock-period mask}: log$_{10}$ Benioff over days 0--25
only (the last 5-day sub-window dropped). Macro AUC collapses to 0.530.

A formal paired-bootstrap test on shared resamples (B = 1000) gives
$\Delta\mathrm{AUC}$ (main $-$ shift) $= +0.267$, $\mathrm{CI}_{95}$
$[+0.20, +0.32]$, with the CI excluding 0. The signal is exclusively
foreshock-period-localized: there is no $> 5$-day-ahead precursor
content recoverable from this scalar.

\subsection{Quantitative magnitude scaling}
\label{sec:results-magnitude}

Stratifying the 139 precursor windows by target magnitude under the
trivial scalar yields:

\begin{center}
\begin{tabular}{lcrcc}
\toprule
Bin & $n_\mathrm{pre}$ & Macro AUC & $\mathrm{CI}_{95}$ \\
\midrule
$M \in [4.5, 4.8)$ & 72 & 0.701 & [0.61, 0.82] \\
$M \in [4.8, 5.2)$ & 33 & 0.837 & [0.76, 0.92] \\
$M \ge 5.2$        & 34 & 0.893 & [0.81, 0.96] \\
\midrule
$M \in [4.5, 5.0)$ & 91 & 0.731 & [0.64, 0.83] \\
$M \ge 5.0$        & 48 & 0.879 & [0.81, 0.93] \\
\bottomrule
\end{tabular}
\end{center}

Per-region AUCs in the $M \ge 5.0$ bin are uniformly elevated:
California $0.884$ ($\mathrm{CI}_{95}\,[0.77, 0.98]$, $n=16$), Cascadia
$0.940$ ($[0.86, 1.00]$, $n=13$), Turkey $0.777$ ($[0.55, 0.94]$,
$n=8$), Italy $0.914$ ($[0.77, 1.00]$, $n=11$). The 3-bin Mann-Kendall
trend test gives $\tau = +1.00$ but is not statistically significant
($p = 0.33$) because three points cannot reach significance under
Mann-Kendall regardless of monotonicity. We therefore report the
trend qualitatively: macro AUC scales monotonically with target $M$
across both 2-bin and 3-bin partitions, with each successive bin's
$\mathrm{CI}_{95}$ overlapping the next but the central tendency
strictly increasing.

This monotonic scaling is consistent with the Gutenberg-Richter
scaling of triggered seismicity (larger mainshocks have proportionally
more and larger preceding events on average); we deliberately do not
attribute this to ETAS-cascade structure specifically, since
distinguishing genuine ETAS triggering from a passive GR effect
would require an ETAS simulation with matched $M_c$, declustering, and
window definition that is beyond this study's scope.

\subsection{TLS template scan: ablation showing it underperforms}
\label{sec:results-ablation}

The TLS-style template-correlation scan of \S\ref{sec:methods-tls},
applied to the same Benioff trajectory, yields macro AUC = $0.704$
(cross-region $\mathrm{CI}_{95}\,[0.64, 0.77]$, $z_\mathrm{iid} =
+8.0$). This passes the pre-reg 3$\sigma$ + Bonferroni gates, but
\emph{underperforms the trivial scalar by 0.075 AUC}, in 4 of 4 LORO
splits (California 0.733 vs 0.811, Cascadia 0.809 vs 0.905, Turkey
0.608 vs 0.691, Italy 0.665 vs 0.710). The Pearson correlation step in
TLS normalizes by within-window standard deviation, discarding the
absolute-magnitude information that the trivial scalar preserves and
that turns out to be the dominant signal carrier. We document this as
a methodological negative result: template-matching machinery does not
add information here; the simplest scalar wins. The
$n$-above-$M_c$ trajectory under the same TLS scan yields macro AUC
$= 0.542$ ($z = +1.62$), sub-3$\sigma$.

% =============================================================================
\section{Discussion}
\label{sec:discussion}

\subsection{What the failure-mode catalogue contributes}

Three of the five modes (\#1 overlap-rule density, \#2 ANSS-as-
teleseismic, \#4 test-region inaccessibility) are operational lessons
familiar individually to subduction-zone-experienced researchers, but
their concrete combination under our pre-registered protocol made them
load-bearing in a way that is worth recording for the next group
attempting cross-regional pre-registered precursor work. Modes \#3
(\(b\)-feature $\times$ $M_c$-window incompatibility) and \#5
(precursor-window-includes-foreshock-period) are the more substantive
contributions: \#3 rules out short-window $b$-drop precursor tests at
moderate $M_c$ as a feasibility class; \#5 highlights an interpretive
gap in the literature definition of ``precursor window'' that our
controls (\S\ref{sec:results-controls}) make sharp.

\subsection{Cross-regional confirmation of foreshock pervasiveness}

The trivial-scalar finding is a quantitative cross-regional
confirmation of the \citet{trugman2019} southern-California finding:
foreshock activity is elevated before most $M{\ge}4.5$ mainshocks, and
the elevation is detectable with the simplest possible scalar at macro
AUC = $0.78$ across four regions on three continents. We add (i)~a
cross-regional AUC quantification with explicit per-region values, (ii)
a dependence-aware significance estimate ($z_\mathrm{block} = +8.4$)
that is more conservative than the i.i.d.\ value, (iii)~a paired
control with formal $\Delta\mathrm{AUC}$ confidence bounds, and (iv)
M-scaling with bootstrap CIs per bin and per region. The
\citet{helmstetter2003a, helmstetter2003b} ETAS framework is consistent
with the monotonic M-scaling we measured, but distinguishing ETAS-
specific cascade structure from generic Gutenberg-Richter triggered-
seismicity scaling is beyond our scope.

\subsection{What the foreshock-period concentration means}

The window-shift, placebo, and mask controls together establish that
the entire signal lives in the final 5 days of the 30-day pre-event
window. This is operationally important: the result is foreshock
detection, not 25-day-ahead long-distance precursor. The pre-reg's
literal definition of ``precursor window'' as $[t-30, t)$ admits the
foreshock period; under the literal protocol, our pass is real, but
under any interpretation that distinguishes foreshocks from long-
distance precursors, the long-distance signal is null.

A reasonable extension would re-run with a precursor window that
explicitly excludes the final 5--7 days, testing whether any sub-
detection-threshold long-distance signal exists. We did not do so
because it would be a post-hoc deviation from pre-reg v1; we flag it
as the natural follow-up and note that under the trivial scalar,
truncating to days 0--25 already produces macro AUC $= 0.530$, leaving
little room for a long-distance signal to be hidden under the same
estimator.

\subsection{Pre-registration discipline: binding vs amended}
\label{sec:disc-prereg}

The PRA-2 amendment was committed before re-evaluation but \emph{after}
v1 surfaced the failure modes. This is the exact degree of freedom that
strict pre-registration aims to eliminate. We are explicit: the v2
protocol is not binding pre-registration in the strictest sense; it is
an amended protocol with documented amendment grounds and full chain
of custody. We mitigated three ways. First, every amendment is
substantively justified by an execution-time failure mode that would
have caused the v1 evaluation to be uninformative regardless of
outcome (mode \#2 forces $b$-features uncomputable; mode \#1 leaves
zero qualifying windows). Second, the amendment was committed to a
public repository before any re-evaluation was run, so the chain
$\mathrm{v1 \to amendment \to re\text{-}evaluation}$ is verifiable.
Third, neither v1 nor v2 had any positive scalar finding; the trivial-
scalar headline was discovered in post-v2 follow-up
(\S\ref{sec:methods-scalar}) and is independent of the v1$\to$v2
amendment in the sense that v1's overlap rule and ComCat catalog do
not affect the trivial scalar's AUC computation conditional on the
qualifying-window set.

A reader who weights only fully-binding pre-registration may discount
the headline finding by treating it as exploratory; we accept this
framing. The five failure modes and the chain-of-custody artifact
remain valuable independently of pre-registration strictness.

\subsection{Score honest, not inflated}

A natural temptation after the trivial-scalar positive is to claim a
``first cross-regional precursor finding.'' We resist this. The
foreshock-pervasiveness phenomenon is well-documented per-region
\citep{trugman2019, lippiello2025, khan2025, moutote2021}; our work
adds a quantitative cross-regional AUC, dependence-aware significance,
controls with formal paired-test, and M-scaling under partial pre-
registration discipline, but does not discover new physics or
methodology beyond the simplest possible scalar. Both the trivial
scalar and the TLS apparatus are deliberately conservative ablations:
neither claims any post-hoc model search beyond the pre-reg-promised
``feature time series'' framing. The honest contribution is the
combination of the failure-mode catalogue, the chain-of-custody
artifact, and the cross-regional AUC at the simplest defensible
estimator.

% =============================================================================
\section{Conclusion}
\label{sec:conclusion}

We pre-registered a CSEP-style cross-regional precursor evaluation and
ran it under the locked protocol with one documented post-v1
amendment. Five execution-time failure modes surfaced, all catalogued
with diagnostic evidence. The pre-registered per-feature scalar and
joint-classifier panels yield clean nulls. The simplest possible
scalar — log$_{10}$ Benioff energy in the final 5 days of the 30-day
pre-event window — passes 3$\sigma$ and Bonferroni at macro AUC $=
0.78$, $z_\mathrm{block} = +8.4$ across four training regions; this
signal is exclusively foreshock-period-localized under three controls
and scales monotonically with target magnitude. A TLS-style template-
correlation apparatus over the same trajectory underperforms the
trivial scalar by 0.075 AUC, documented as a methodological negative
result. We frame the headline as cross-regional confirmation of
\citet{trugman2019} foreshock-pervasiveness with the simplest
defensible estimator, not as a novel discovery. The five-mode
failure catalogue and the publicly verifiable chain of custody are the
methodological contributions.

\section*{Acknowledgements}
The author thanks the ISC for maintaining the global bulletin, the
NCEDC for hosting California waveform data, and INGV/MedNet for the
post-2009 successor station MN.AQU. ObsPy \citep{beyreuther2010} and
\texttt{usgs-libcomcat} provided the data-access infrastructure.
Pre-registration drafting, post-detection control design, and post-
peer-review baseline ablation used Anthropic Claude as an interactive
collaborator (Claude is acknowledged but not a co-author per arXiv
policy).

\section*{Data and code availability}
\begin{sloppypar}
All code, the pre-registration document \citep{loftus2026prereg}, the
amendment \citep{loftus2026pra2}, per-region cached catalogs, per-window
feature CSVs, and reproducible experiment scripts are publicly available
at \citet{loftus2026repo}. Each work session is a single git commit; the
chain of custody from pre-reg through results is verifiable without
access to any private artifact. Specifically, the trivial-scalar
headline (\S\ref{sec:results-scalar-finding}) is reproduced by
\texttt{exp18\_trivial\_baseline\_ablation/run.py}, the controls
(\S\ref{sec:results-controls}) by
\texttt{exp19\_trivial\_scalar\_controls/run.py}, and the M-scaling
(\S\ref{sec:results-magnitude}) by
\texttt{exp20\_trivial\_scalar\_magnitude\_scaling/run.py} (all under
\texttt{experiments/}).
\end{sloppypar}

\bibliography{references}
\end{document}
